% ============================================================================
%  HakuHook — Informe de desarrollo del Gateway Hub
%  Documento vivo: se actualiza conforme avanza el proyecto.
%  Compilar con: pdflatex informe-hakuhook.tex
% ============================================================================
\documentclass[11pt,a4paper]{article}

\usepackage[utf8]{inputenc}
\usepackage[T1]{fontenc}
\usepackage[spanish,es-tabla]{babel}
\usepackage[margin=2.5cm]{geometry}
\usepackage{xcolor}
\usepackage{listings}
\usepackage{booktabs}
\usepackage{array}
\usepackage{enumitem}
\usepackage{parskip}
\usepackage{fancyhdr}
\usepackage[hidelinks]{hyperref}

% ---------- Colores ----------
\definecolor{codebg}{RGB}{248,248,248}
\definecolor{codeframe}{RGB}{210,210,210}
\definecolor{kw}{RGB}{0,80,160}
\definecolor{typekw}{RGB}{140,60,20}
\definecolor{cmt}{RGB}{90,120,90}
\definecolor{str}{RGB}{170,60,60}

% ---------- Lenguaje TypeScript para listings ----------
\lstdefinelanguage{TypeScript}{
  keywords={const,let,var,function,return,if,else,throw,new,for,of,in,import,from,export,type,interface,class,extends,implements,void,true,false,null},
  keywordstyle=\color{kw}\bfseries,
  ndkeywords={Map,Set,WebSocket,ClientInfo,Room,PendingApproval,JsonRpcResponse,JsonRpcError,ClientRole},
  ndkeywordstyle=\color{typekw}\bfseries,
  sensitive=false,
  comment=[l]{//},
  morecomment=[s]{/*}{*/},
  morestring=[b]',
  morestring=[b]"
}

\lstset{
  language=TypeScript,
  basicstyle=\ttfamily\small,
  backgroundcolor=\color{codebg},
  frame=single,
  rulecolor=\color{codeframe},
  breaklines=true,
  columns=fullflexible,
  showstringspaces=false,
  commentstyle=\color{cmt}\itshape,
  stringstyle=\color{str},
  tabsize=2,
  captionpos=b,
  aboveskip=0.8em,
  belowskip=0.8em
}

% ---------- Encabezado / pie ----------
\pagestyle{fancy}
\fancyhf{}
\fancyhead[L]{HakuHook --- Informe de desarrollo}
\fancyhead[R]{gateway-hub}
\fancyfoot[C]{\thepage}

\title{\textbf{HakuHook --- Informe de desarrollo}\\ \large Gateway Hub: WebSocket relay HITL}
\author{Haku}
\date{\today}

\begin{document}

\maketitle
\tableofcontents
\newpage

% ===========================================================================
\section{Introducción y visión del producto}
% ===========================================================================

HakuHook es un sistema de \textbf{guardrails y arbitraje multi-agente en tiempo
real} (\textbf{HITL}, \emph{human in the loop}) cuyo objetivo es prevenir el
\emph{scope creep} y la degradación arquitectónica cuando desarrolladores junior
usan agentes de IA (Antigravity, Cursor, Claude Code) en proyectos corporativos.

El flujo general es:

\begin{enumerate}
  \item Un hook transparente intercepta \textbf{exclusivamente} tool calls
        destructivos o de alto impacto (\texttt{write\_file},
        \texttt{create\_file}, \texttt{execute\_command}, migraciones de base de datos).
  \item Retiene la promesa de ejecución y envía el contexto
        (prompt del junior + tool call + diff propuesto) al entorno del
        \textbf{Tech Lead / Senior}.
  \item En la máquina del senior, un \textbf{LLM evaluador} integrado con
        Gentle-AI / RDD / skills arquitectónicas analiza el impacto y presenta
        un veredicto estructurado: \texttt{[Aprobar]},
        \texttt{[Rechazar con feedback pedagógico]}, \texttt{[Modificar]}.
\end{enumerate}

El modelo de roles es asimétrico por diseño:

\begin{itemize}
  \item \textbf{Junior}: el agente de IA que \emph{propone} y \emph{pide permiso}
        para ejecutar acciones destructivas.
  \item \textbf{Senior}: el humano (HITL) que \emph{arbitra}: revisa el contexto
        y \emph{decide} aprobar, rechazar o modificar.
\end{itemize}

\section{Arquitectura del monorepo}

Monorepo TypeScript con pnpm workspaces bajo el scope \texttt{@hakuhook/*}.

\begin{table}[h]
  \centering
  \begin{tabular}{@{}p{3.2cm}p{7.5cm}p{2.8cm}@{}}
    \toprule
    \textbf{Paquete} & \textbf{Rol} & \textbf{Estado} \\
    \midrule
    \texttt{shared} & Tipos comunes, interfaces de sesión,
                      contratos JSON-RPC 2.0 & Completado \\
    \texttt{gateway-hub} & Servidor WebSocket relay; salas por \texttt{projectId};
                           sincroniza peticiones junior $\leftrightarrow$ senior
                           en tiempo real & \textbf{En desarrollo} \\
    \texttt{proxy-mcp} & Interceptor STDIO para MCP; filtra llamadas y se
                         comunica vía WebSocket con el hub & Pendiente \\
    \texttt{evaluator} & Motor de arbitraje del senior; evalúa diff con LLM;
                         UI interactiva en terminal & Pendiente \\
    \texttt{cli} & Binario global (\texttt{hakuhook init}, \texttt{hakuhook login});
                   aislamiento por repositorio & Pendiente \\
    \bottomrule
  \end{tabular}
  \caption{Paquetes del monorepo y su estado.}
\end{table}

\section{Estado del roadmap}

\begin{table}[h]
  \centering
  \begin{tabular}{@{}p{1.6cm}p{5.2cm}p{3cm}p{3.4cm}@{}}
    \toprule
    \textbf{Paso} & \textbf{Contenido} & \textbf{Estado} & \textbf{Evidencia} \\
    \midrule
    1 & Tipos base (\texttt{types.ts}) & Completado & Commit \texttt{2c8054a} \\
    2 & Utilidades JSON-RPC (\texttt{json-rpc.ts}) & Completado & Commit \texttt{2c8054a} \\
    3 & Registro de salas (\texttt{room-manager.ts}) & Completado & Commits \texttt{93bf9a5}, \texttt{561a96e} \\
    4 & Handlers (routing junior $\rightarrow$ senior) & En desarrollo &
           \texttt{register-session.ts} completo; \texttt{request-approval.ts}
           con firma, cuerpo pendiente \\
    5 & Carnet de conexiones (\texttt{connection-store.ts}) & Completado & Sesión 2026-09-02 \\
    6 & Servidor WebSocket + dispatcher (\texttt{server.ts}) & Pendiente & --- \\
    7 & Punto de entrada (\texttt{index.ts}) + build & Pendiente & --- \\
    \bottomrule
  \end{tabular}
  \caption{Roadmap de implementación del gateway-hub.}
\end{table}

El build actual pasa correctamente: \texttt{pnpm build} (TypeScript compila sin errores).

\section{Detalle paso a paso}

\subsection{Paso 1: Tipos base}

Define la ``tarjeta'' de cada participante, la sala y las peticiones pendientes.

\begin{lstlisting}[caption=packages/gateway-hub/src/types.ts]
import type { WebSocket } from "ws";
import type { ClientRole } from "@hakuhook/shared";

export interface ClientInfo {
  socket: WebSocket;
  projectId: string;
  role: ClientRole;
}

export interface Room {
  projectId: string;
  senior: ClientInfo | null;
  juniors: Set<ClientInfo>
}

export interface PendingApproval {
  projectId: string;
  requester: ClientInfo;
}
\end{lstlisting}

Tipos compartidos relevantes (definen el contrato de negocio):

\begin{lstlisting}[caption=packages/shared/src/types.ts (extracto)]
export type ClientRole = "junior" | "senior";

export interface ApprovalDecision {
  requestId: string | number;
  status: "approved" | "rejected" | "modified";
  feedback?: string;
  modifiedArguments?: Record<string, any>;
}
\end{lstlisting}

\begin{itemize}
  \item \texttt{ClientInfo} es la tarjeta: conexión, sala a la que pertenece y rol.
  \item \texttt{Room} modela la asimetría: una casilla única de senior
        (\texttt{ClientInfo | null}) y una colección de juniors (\texttt{Set}).
  \item \texttt{PendingApproval} no lleva \texttt{id}: la llave (el
        \texttt{approvalId}) se genera en los handlers (Paso 4).
\end{itemize}

\subsection{Paso 2: Utilidades JSON-RPC}

Helpers para enviar respuestas y errores por el socket según el contrato JSON-RPC 2.0.

\begin{lstlisting}[caption=packages/gateway-hub/src/json-rpc.ts]
import type { WebSocket } from "ws";
import type { JsonRpcError, JsonRpcResponse } from "@hakuhook/shared";

export function send(socket: WebSocket, message: unknown): void {
  socket.send(JSON.stringify(message))
}

export function sendResult(
  socket: WebSocket,
  id: string | number,
  result: unknown,
): void {
  const response: JsonRpcResponse = { jsonrpc: "2.0", id, result };
  send(socket, response);
}

export function sendError(
  socket: WebSocket,
  id: string | number,
  code: number,
  message: string,
): void {
  const error: JsonRpcError = { code, message };
  const response: JsonRpcResponse = { jsonrpc: "2.0", id, error };
  send(socket, response);
}
\end{lstlisting}

\subsection{Paso 3: Room manager}

El \texttt{room-manager} es un \textbf{registro pasivo}: sabe quién está en cada
sala, pero no envía mensajes. Los handlers (Paso 4) lo consultan para enrutar.

Diseño: módulo-singleton con dos contenedores a nivel de módulo:

\begin{lstlisting}[caption=packages/gateway-hub/src/room-manager.ts (vista corregida)]
import type { ClientInfo, PendingApproval, Room } from "./types.js";

const rooms = new Map<string, Room>();
const approvals = new Map<string, PendingApproval>()

export function joinRoom(client: ClientInfo): Room {
  let room = rooms.get(client.projectId)

  if (client.role == "senior" && room?.senior) {
    throw new Error(
      `Project "${client.projectId}" already has an active senior`,
    )
  }

  if (!room) {
    room = {
      projectId: client.projectId,
      senior: null,
      juniors: new Set(),
    }
    rooms.set(client.projectId, room)
  }

  if (client.role == "senior") {
    room.senior = client
  } else {
    room.juniors.add(client)
  }

  return room
}

export function removeClient(client: ClientInfo): void {
  const room = rooms.get(client.projectId)
  if (!room) return

  // Capa 1: quitar de la sala segun su contenedor
  if (client.role == "senior") {
    if (room.senior === client) {
      room.senior = null
    }
  } else {
    room.juniors.delete(client)
  }

  // Capa 2: purgar aprobaciones fantasma (aplica a cualquier rol)
  for (const [approvalId, approval] of approvals) {
    if (approval.requester === client) {
      approvals.delete(approvalId)
    }
  }

  // Capa 3: GC de sala vacia (aplica a cualquier rol)
  if (!room.senior && room.juniors.size === 0) {
    rooms.delete(client.projectId)
  }
}

export function getSenior(projectId: string): ClientInfo | null {
  const senior = rooms.get(projectId)?.senior
  if (!senior) return null
  return { ...senior }
}

export function getJuniors(projectId: string): ClientInfo[] {
  const room = rooms.get(projectId)
  if (!room) return []
  return [...room.juniors].map((junior) => ({ ...junior }))
}

export function requestApproval(approvalId: string, client: ClientInfo): PendingApproval {
  const approval: PendingApproval = {
    projectId: client.projectId,
    requester: client,
  }
  approvals.set(approvalId, approval)
  return approval
}

export function takeApproval(approvalId: string): PendingApproval | null {
  const approval = approvals.get(approvalId)
  if (!approval) return null
  approvals.delete(approvalId)
  return approval
}
\end{lstlisting}

\textbf{Nota:} la versión mostrada incluye la corrección documentada en la
Sección~\ref{sec:hallazgos}. El código \emph{commiteado} en el repositorio aún
tiene el error pendiente de aplicar.

\subsubsection{Las tres capas de \texttt{removeClient}}

\begin{itemize}
  \item \textbf{Capa 1 (detach):} quitar la tarjeta según su contenedor.
        El junior se borra de la colección (\texttt{Set.delete}); la casilla del
        senior se vacía asignando \texttt{null}.
  \item \textbf{Capa 2 (fantasmas):} las peticiones pendientes del que se va se
        eliminan de la bandeja. La llave de \texttt{approvals} es el
        \texttt{approvalId} (no la persona), por lo que hay que barrer todas
        las entradas con un \texttt{for} y comparar por identidad.
  \item \textbf{Capa 3 (GC):} si la sala quedó sin senior y sin juniors, se
        elimina del registro. La cierra \emph{el último que sale},
        independientemente del rol.
\end{itemize}

\subsubsection{Mecanismos de almacenamiento}

Cada contenedor se vacía con su herramienta correcta:

\begin{table}[h]
  \centering
  \begin{tabular}{@{}p{3.6cm}p{4.4cm}p{5cm}@{}}
    \toprule
    \textbf{Contenedor} & \textbf{Qué es} & \textbf{Cómo se vacía} \\
    \midrule
    \texttt{room.senior} & Casilla única (propiedad) & \texttt{room.senior = null} \\
    \texttt{room.juniors} & Colección (\texttt{Set}) & \texttt{room.juniors.delete(client)} \\
    \texttt{rooms} / \texttt{approvals} & Cajonera (\texttt{Map}) & \texttt{<map>.delete(llave)} \\
    \bottomrule
  \end{tabular}
  \caption{Formas de quitar elementos según el contenedor.}
\end{table}

\subsubsection{Asimetría senior vs.\ junior}

\begin{table}[h]
  \centering
  \begin{tabular}{@{}p{3.2cm}p{5.4cm}p{5.4cm}@{}}
    \toprule
    & \textbf{Junior} & \textbf{Senior} \\
    \midrule
    Qué es & Agente de IA del dev & Tech Lead humano (HITL) \\
    Qué hace & Pide permiso (\texttt{requestApproval}, es el \texttt{requester}) &
               Toma el permiso (\texttt{takeApproval}) y decide
               (\texttt{ApprovalDecision}) \\
    Cantidad por sala & Muchos (\texttt{Set}) & Uno (casilla única) \\
    Regla especial & --- & Máximo un senior por sala (lanza error) \\
    \bottomrule
  \end{tabular}
  \caption{El junior propone; el senior arbitra.}
\end{table}

\subsection{Paso 4: Handlers -- el routing junior $\rightarrow$ senior}

En este paso el registro ``cobra vida'': los clientes, una vez en su sala,
empiezan a \emph{pedir permiso}. Los handlers reciben el mensaje JSON-RPC ya
enrutado y lo convierten en acciones sobre las salas y las peticiones.

\subsubsection{El handler de registro (\texttt{register-session.ts})}

Su trabajo es \textbf{asociar una conexión a una sala y a un rol}. Es el
\emph{único} lugar donde el servidor aprende ``quién es'' cada socket.

\begin{lstlisting}[caption=packages/gateway-hub/src/handlers/register-session.ts]
import type { JsonRpcRequest, RegisterSessionParams } from "@hakuhook/shared";
import type { WebSocket } from "ws";
import type { ClientInfo } from "../types.js";
import { joinRoom } from "../room-manager.js"
import { sendError, sendResult } from "../json-rpc.js";

export function handlerSession(
  jsonRpc: JsonRpcRequest<RegisterSessionParams>,
  socket: WebSocket
): void {
  const { projectId, role } = jsonRpc.params;

  if (role !== "junior" && role !== "senior") {
    sendError(socket, jsonRpc.id!, -32602, `Invalid role: "${role}"`);
    return;
  }

  const client: ClientInfo = { socket, projectId, role };

  try {
    const room = joinRoom(client);
    sendResult(socket, jsonRpc.id!, { projectId: room.projectId, role });
  } catch (err) {
    const message = err instanceof Error ? err.message : "Unknown error";
    sendError(socket, jsonRpc.id!, -32000, message)
  }
}
\end{lstlisting}

Puntos clave de su lógica:

\begin{itemize}
  \item \textbf{Generic \texttt{RegisterSessionParams}}: le dice a TypeScript que
        los \texttt{params} del mensaje traen \texttt{projectId} y \texttt{role};
        sin el generic, \texttt{jsonRpc.params} sería \texttt{unknown} y no se
        podría leer \texttt{params.projectId}.
  \item \textbf{Validar antes de mutar}: se valida el \texttt{role} con
        \texttt{sendError(-32602)} \emph{antes} de tocar el estado. El error
        JSON-RPC \texttt{-32602} es ``Invalid params''.
  \item \textbf{\texttt{try/catch}:} el único riesgo real es \texttt{joinRoom}
        (que lanza si la sala ya tiene senior). Si falla, se responde
        \texttt{sendError(-32000)} (error de negocio/servidor), y el cambio no
        se aplica.
  \item \textbf{Non-null assertion \texttt{id!}}: \texttt{id} es opcional en el
        contrato, pero acá se necesita sí o sí para responder; el \texttt{!}
        ``confirma'' a TypeScript que existe.
\end{itemize}

\textbf{Pendiente de conectar:} este handler aún \emph{no} llama a
\texttt{registerClient} cuando construye el \texttt{ClientInfo}; sin ese cable,
el carnet de conexiones nunca se llena y todo \texttt{request-approval}
devolvería ``Not registered''.

\subsubsection{El handler de petición (\texttt{request-approval.ts})}

Por ahora solo tiene la \textbf{firma} escrita; el cuerpo está vacío. La firma
ya marca una decisión clave de diseño:

\begin{lstlisting}[caption=packages/gateway-hub/src/handlers/request-approval.ts (esqueleto)]
import type { JsonRpcRequest, DestructiveToolParams } from "@hakuhook/shared";
import type { WebSocket } from "ws";

export function handlerRequestApproval(
  jsonRpc: JsonRpcRequest<DestructiveToolParams>,
  socket: WebSocket
): void {

}
\end{lstlisting}

\textbf{¿Para qué se tipa con \texttt{DestructiveToolParams}?} Para leer el
\emph{contenido} de la petición (\texttt{toolName}, \texttt{arguments},
\texttt{diff}, \texttt{prompt}). Es importante notar que este tipo \textbf{no
trae \texttt{projectId}}: el \texttt{projectId} se obtiene del carnet de
conexiones mediante \texttt{getClient(socket)}, nunca de los \texttt{params}.

\subsubsection{El carnet de conexiones (\texttt{connection-store.ts})}

Nuevo módulo que resuelve la pregunta \emph{``¿quién es el dueño de este
socket?''}. Es el ``carnet de identidad por conexión''.

\begin{lstlisting}[caption=packages/gateway-hub/src/connection-store.ts]
import type { WebSocket } from "ws";
import type { ClientInfo } from "./types.js";

const clients = new Map<WebSocket, ClientInfo>();

export function registerClient(client: ClientInfo): void {
  clients.set(client.socket, client);
}

export function getClient(socket: WebSocket): ClientInfo | undefined {
  return clients.get(socket);
}

export function removeClientConnection(socket: WebSocket): void {
  clients.delete(socket);
}
\end{lstlisting}

Puntos clave:

\begin{itemize}
  \item \textbf{Map privado por conexión}: \texttt{Map<WebSocket, ClientInfo>}
        relaciona cada socket con su tarjeta (que incluye el \texttt{projectId}).
        Es un carnet distinto del \texttt{rooms} (que es por \emph{sala}).
  \item \textbf{Guardianía}: \texttt{clients} es privado; se expone solo
        \texttt{registerClient}, \texttt{getClient} y
        \texttt{removeClientConnection}. Nadie muta el carnet desde afuera; del
        mismo modo que \texttt{getSenior} devuelve copias.
  \item \textbf{Responsabilidad única}: se separa del \texttt{server.ts} para que
        el servidor solo \emph{orqueste} y no acumule estado (evita un archivo
        ``que hace de todo'').
\end{itemize}

\subsubsection{Separación de responsabilidades en el Paso 4}

\begin{table}[h]
  \centering
  \begin{tabular}{@{}p{4cm}p{5cm}p{4.4cm}@{}}
    \toprule
    \textbf{Módulo} & \textbf{Posee} & \textbf{Responde a} \\
    \midrule
    \texttt{connection-store.ts} & Map \texttt{clients} (socket $\rightarrow$ ClientInfo) & ¿Quién es este socket? \\
    \texttt{room-manager.ts} & Map \texttt{rooms} (projectId $\rightarrow$ Room) & ¿Quién es el senior/junior de esta sala? \\
    \texttt{server.ts} (futuro) & Flujo: wss + \texttt{message} + dispatcher & ¿A qué handler enrutar? \\
    \bottomrule
  \end{tabular}
  \caption{Dos carnets distintos, dos dueños distintos.}
\end{table}

\section{Conceptos clave aprendidos}

\begin{description}
  \item[Map y sus métodos] Un \texttt{Map} es una cajonera: \texttt{get(llave)}
        abre un cajón (o devuelve \texttt{undefined}), \texttt{set} guarda,
        \texttt{delete} quita entrada por llave.
  \item[Tipos genéricos] \texttt{Map<string, Room>} se lee ``cajonera con
        etiquetas string que guarda Rooms'': primer parámetro = tipo de llave,
        segundo = tipo de valor.
  \item[Identidad por referencia] \texttt{===} entre objetos compara la
        dirección de memoria, no el contenido. Dos tarjetas que se ven iguales
        son dos objetos distintos. Es la base de \texttt{Set.delete} y de
        \texttt{room.senior === client}.
  \item[Casilla única vs.\ colección] El senior vive en una casilla
        (\texttt{ClientInfo | null}); los juniors en una colección. Por eso se
        vacían distinto.
  \item[Checkout semántico] \texttt{takeApproval} \emph{saca} la petición de la
        bandeja: al tomarla, nadie más puede tomarla. Tomar = sacar.
  \item[Validar antes de mutar] Comprobar invariantes (ej.\ un solo senior)
        antes de modificar el estado.
  \item[Lazy init] Crear la sala solo cuando llega la primera persona; no
        pre-crear nada.
  \item[GC de sala vacía] Borrar del registro la sala que se queda sin nadie
        (el último apaga la luz).
  \item[Snapshot copy] \texttt{getSenior}/\texttt{getJuniors} devuelven copias
        (\texttt{\{ ...x \}}), no referencias internas: se puede leer, no mutar.
  \item[Type assertion (\texttt{as})] \texttt{JSON.parse} devuelve \texttt{unknown}.
        \texttt{parsed as JsonRpcRequest} es una ``confianza'' que le damos a
        TypeScript para poder leer \texttt{.method}/\texttt{.params}; \textbf{no
        valida en runtime}.
  \item[El socket = identidad + destino] Cada cliente tiene UN socket: es el
        canal para responderle (\texttt{socket.send}) y la llave para saber quién
        es (\texttt{getClient(socket)}). Por eso ``siempre'' aparece.
  \item[El dispatcher] ``Recepcionista'': un único punto de entrada que lee
        \texttt{jsonRpc.method} y enruta a cada handler con un \texttt{switch}.
  \item[El generic tipa el contenido, no el proyecto] \texttt{JsonRpcRequest<T>}
        sirve para leer los \texttt{params} del mensaje; el \texttt{projectId}
        del routing sale del carnet (\texttt{getClient}), no de los params.
  \item[Guardianía del estado] El dueño de un Map expone funciones
        (\texttt{registerClient}/\texttt{getClient}) en lugar de dejar tocar el
        Map directo; del mismo modo que \texttt{getSenior} vuelve copias.
\end{description}

\section{Resumen de los ``porqués'' (para no perder la sintonía)}

Resumen comprimido de las decisiones y su razonamiento, para retomar el hilo
rápido en la próxima sesión.

\subsection{Por qué un \texttt{Map} de \texttt{clients}}
Para saber \emph{quién es} cada socket. No basta con guardar sockets sueltos:
se guarda la relación \texttt{socket $\rightarrow$ ClientInfo} (que trae
\texttt{projectId} y \texttt{role}). Responde a ``¿de qué junior viene este
mensaje?''.

\subsection{Por qué \texttt{server.ts} no guarda el carnet; \texttt{connection-store} sí}
Para no crear un ``archivo que hace de todo''. \texttt{server.ts} solo
orquesta (crea el wss, escucha, enruta); \texttt{connection-store.ts} es el
dueño exclusivo del carnet de conexiones. Una responsabilidad por módulo.

\subsection{Por qué \texttt{getClient} y no ``reusar'' \texttt{getSenior}}
Son carnets \textbf{distintos}:
\begin{itemize}
  \item \texttt{getSenior}/\texttt{getJuniors} le preguntan al Map
        \texttt{rooms} (por \emph{sala}/\texttt{projectId}).
  \item \texttt{getClient} le pregunta al Map \texttt{clients} (por
        \emph{conexión}/\texttt{socket}).
\end{itemize}
Cuando llega un \texttt{request-approval}, lo \'unico que se tiene es el
socket; \texttt{getClient} es el puente para obtener el \texttt{projectId} y,
\emph{después}, poder usar \texttt{getSenior(projectId)}.

\subsection{Por qué tipar \texttt{DestructiveToolParams} si no tiene \texttt{projectId}}
El generic tipa el \emph{contenido} del mensaje (\texttt{toolName},
\texttt{arguments}, \texttt{diff}, \texttt{prompt}). El \texttt{projectId} sale
\textbf{por separado} del carnet (\texttt{getClient(socket)}), no de los
params. Son dos canales de información distintos.

\subsection{Por qué el \texttt{as JsonRpcRequest} en \texttt{dispatch}}
\texttt{JSON.parse} devuelve \texttt{unknown}; \texttt{as} es una afirmación de
tipo para leer \texttt{.method} y \texttt{.params}. Es compile-time, no valida
en runtime; por eso se combina con el \texttt{try/catch} del parse.

\subsection{Por qué \texttt{dispatch} necesita \texttt{jsonRpc}}
Para leer \texttt{jsonRpc.method} (decide a qué handler va) y para entregarle
ese mismo mensaje al handler (que lee sus \texttt{params}). Es el sobre con la
carta: el recepcionista mira el remitente para clasificar y entrega el mismo
sobre a la oficina.

\subsection{Por qué ``siempre'' aparece el \texttt{socket}}
Es a la vez el \textbf{destino} (a quién responder: \texttt{sendError(socket)})
y la \textbf{identidad} (de quién viene: \texttt{getClient(socket)}). Es el
único vínculo entre el servidor y cada cliente, por eso viaja en todo el flujo.

\section{Hallazgos y correcciones}
\label{sec:hallazgos}

\begin{table}[h]
  \centering
  \begin{tabular}{@{}p{2.4cm}p{6.4cm}p{2.4cm}p{2.6cm}@{}}
    \toprule
    \textbf{Fecha} & \textbf{Hallazgo} & \textbf{Impacto} & \textbf{Estado} \\
    \midrule
    2026-08-30 & \texttt{removeClient}: toda la limpieza quedó anidada dentro
                 de \texttt{if (client.role == "senior")}. Consecuencias:
                 (1) un junior que se va NO se elimina de la sala;
                 (2) sus aprobaciones fantasma no se purgan;
                 (3) la sala nunca se cierra cuando el último en irse es junior
                 (no hay GC). & El registro queda con juniors ``fantasma'' y el
                 relay podría enviar mensajes a sockets muertos. &
                 Corregido en el informe, pendiente de aplicar en el repo. \\
    \bottomrule
  \end{tabular}
  \caption{Errores detectados durante el desarrollo.}
\end{table}

La corrección consiste en sacar la limpieza de aprobaciones y el GC \emph{fuerza}
del bloque \texttt{if (client.role == "senior")}, y dejar el
\texttt{if/else} de rol únicamente para decidir \emph{cómo} se desengancha la
tarjeta (casilla vs.\ colección). Aplicar y volver a compilar:
\texttt{pnpm build}.

\section{Próximos pasos}

\begin{enumerate}
  \item \textbf{Terminar el handler \texttt{request-approval.ts}:} llenar el cuerpo
        usando \texttt{requestApproval} (guardar petición pendiente),
        \texttt{getSenior} (ubicar al senior de la sala) y enviarle la petición
        por su socket; si no hay senior, responder \texttt{sendError}.
  \item \textbf{Conectar \texttt{register-session.ts}:} que al construir el
        \texttt{ClientInfo} llame a \texttt{registerClient} para llenar el
        carnet de conexiones (sin esto, cualquier \texttt{request-approval}
        devolvería ``Not registered'').
  \item \textbf{Crear el servidor WebSocket (\texttt{server.ts}):} instancia el
        \texttt{WebSocketServer}, escucha \texttt{connection}/\texttt{message}/\texttt{close},
        y enruta por \texttt{jsonRpc.method} con un \texttt{switch} (el
        \emph{dispatcher}) hacia los handlers. Uso de \texttt{getClient} para
        obtener el \texttt{projectId} del remitente en \texttt{request-approval}.
  \item \textbf{Punto de entrada:} \texttt{index.ts} + build final.
\end{enumerate}

\section{Bitácora de cambios}

Registro de avances. Agregar una fila por cada cambio significativo.

\begin{table}[h]
  \centering
  \begin{tabular}{@{}p{2.6cm}p{3.2cm}p{7.6cm}@{}}
    \toprule
    \textbf{Fecha} & \textbf{Commit} & \textbf{Descripción} \\
    \midrule
    2026-09-02 & --- & \texttt{connection-store.ts} creado; \texttt{handlers}:
                  \texttt{register-session.ts} completo y
                  \texttt{request-approval.ts} con firma. Diseño del
                  \texttt{server.ts} (wss + dispatcher) acordado.
    2026-08-23 & \texttt{561a96e} & \texttt{removeClient}, \texttt{getSenior},
                  \texttt{getJuniors}, \texttt{requestApproval},
                  \texttt{takeApproval} \\
    2026-08-23 & \texttt{93bf9a5} & \texttt{joinRoom} (validar antes de mutar,
                  lazy init, un solo senior) \\
    2026-08-20 & \texttt{2c8054a} & Tipos base + utilidades JSON-RPC \\
    2026-08-20 & \texttt{f75e348} & Ignorar dependencias (\texttt{.gitignore}) \\
    2026-08-20 & \texttt{7e8aefe} & Tipos del gateway-hub \\
    \bottomrule
  \end{tabular}
  \caption{Bitácora de cambios del gateway-hub.}
\end{table}

\end{document}